
\chapter{Levantamento e Especificação de Requisitos}
\label{chap:requisitos}

% Apresente o objetivo do levantamento de requisitos
% e como os requisitos foram identificados e organizados.

\section{Escopo do sistema}
\label{sec:escopo}

% Descreva a finalidade do DoseCerta, seu público-alvo,
% suas funcionalidades incluídas e os limites do sistema.

\section{Atores do sistema}
\label{sec:atores}

% Identifique os usuários e sistemas externos que interagem
% com o DoseCerta. Descreva suas responsabilidades.

\section{Requisitos funcionais}
\label{sec:requisitos_funcionais}

% Requisitos funcionais descrevem comportamentos
% e funcionalidades que o sistema deve oferecer.

\begin{table}[htbp]
    \centering
    \caption{Requisitos funcionais do DoseCerta}
    \label{tab:requisitos_funcionais}
    \begin{tabular}{|p{2cm}|p{3cm}|p{8cm}|}
        \hline
        \textbf{Identificador} &
        \textbf{Requisito} &
        \textbf{Descrição} \\
        \hline
        RF01 & [Nome] & [Descrição do comportamento esperado.] \\
        \hline
        RF02 & [Nome] & [Descrição do comportamento esperado.] \\
        \hline
        RF03 & [Nome] & [Descrição do comportamento esperado.] \\
        \hline
    \end{tabular}
\end{table}

\section{Requisitos não funcionais}
\label{sec:requisitos_nao_funcionais}

% Requisitos não funcionais descrevem atributos de qualidade,
% restrições e condições de operação do sistema.

\begin{table}[htbp]
    \centering
    \caption{Requisitos não funcionais do DoseCerta}
    \label{tab:requisitos_nao_funcionais}
    \begin{tabular}{|p{2cm}|p{3cm}|p{8cm}|}
        \hline
        \textbf{Identificador} &
        \textbf{Categoria} &
        \textbf{Descrição} \\
        \hline
        RNF01 & Usabilidade & [Condição de usabilidade a ser atendida.] \\
        \hline
        RNF02 & Segurança & [Condição de segurança a ser atendida.] \\
        \hline
        RNF03 & Confiabilidade & [Condição de confiabilidade a ser atendida.] \\
        \hline
    \end{tabular}
\end{table}

\section{Regras de negócio}
\label{sec:regras_negocio}

% Documente regras que restringem ou orientam
% o funcionamento do sistema.

\begin{itemize}
    \item RN01: [Descrição da regra de negócio.]
    \item RN02: [Descrição da regra de negócio.]
\end{itemize}

\section{Restrições do sistema}
\label{sec:restricoes}

% Registre restrições técnicas, físicas, operacionais
% ou relacionadas ao escopo do projeto.

\section{Relação entre requisitos e testes}
\label{sec:requisitos_testes}

% Explique como os requisitos serão verificados.
% Se necessário, apresente uma matriz relacionando
% identificadores de requisitos aos casos de teste.
